Un institut ha preguntat als 300 alumnes de batxillerat quina modalitat fan i si pensen estudiar una
carrera científica ($CC$):
\begin{center}
\begin{tabular}{l|ccc|c}
 & Ciències & Humanitats & Arts & Total\\ \hline
Carrera científica & 90 & 20 & 10 & 120\\
No & 60 & 80 & 40 & 180\\ \hline
Total & 150 & 100 & 50 & 300
\end{tabular}
\end{center}
Es tria un alumne a l'atzar.

\begin{apartats}

\apartat[1,25]{0,75}
Quina és la probabilitat que faci Ciències? I que faci Humanitats o que pensi estudiar una carrera
científica?

\begin{solucio}
$P(\text{Ciències})=\dfrac{150}{300}=0{,}5$.\\
$P(\text{Humanitats}\cup CC)=\dfrac{100+120-20}{300}=\dfrac{200}{300}=\dfrac23\approx0{,}667$: els 20
alumnes d'Humanitats que pensen estudiar una carrera científica no s'han de comptar dues vegades.
\end{solucio}

\begin{tria}{taula-tasca}
\itemtria{original}{1}{1,25}
Si pensa estudiar una carrera científica, quina és la probabilitat que faci Arts? I si fa Ciències, quina
és la probabilitat que no pensi estudiar-ne cap?

\begin{solucio}
$P(\text{Arts}\mid CC)=\dfrac{10}{120}=\dfrac1{12}\approx0{,}083$.\\
$P\left(\overline{CC}\mid\text{Ciències}\right)=\dfrac{60}{150}=0{,}4$.
\end{solucio}

\itemtria{taula-i-bayes}{1}{1,25}
Comprova amb la taula que $P(CC\mid\text{Ciències})=0{,}6$ i que
$P\left(CC\mid\overline{\text{Ciències}}\right)=0{,}2$. Calcula després $P(\text{Ciències}\mid CC)$ amb el
teorema de Bayes, a partir només de $P(\text{Ciències})$ i d'aquestes dues probabilitats, i comprova que
coincideix amb el que dona la taula.

\begin{solucio}
$P(CC\mid\text{Ciències})=\dfrac{90}{150}=0{,}6$ i
$P\left(CC\mid\overline{\text{Ciències}}\right)=\dfrac{20+10}{150}=0{,}2$.\\
Bayes: $P(\text{Ciències}\mid CC)=\dfrac{0{,}5\cdot0{,}6}{0{,}5\cdot0{,}6+0{,}5\cdot0{,}2}
=\dfrac{0{,}3}{0{,}4}=0{,}75$.\\
A la taula, $\dfrac{90}{120}=0{,}75$. Coincideixen: el teorema de Bayes fa el mateix que llegir la taula a
la fila de $CC$, dividir els 90 alumnes de Ciències que la volen estudiar pels 120 que la volen estudiar
en total.
\end{solucio}
\end{tria}

\begin{nomesllarg}
\apartat{0,75}
Són independents «fer Arts» i «pensar estudiar una carrera científica»?

\begin{solucio}
$P(\text{Arts}\cap CC)=\dfrac{10}{300}\approx0{,}033$, però
$P(\text{Arts})\cdot P(CC)=\dfrac16\cdot0{,}4\approx0{,}067$: \textbf{no} són independents.
\end{solucio}
\end{nomesllarg}

\end{apartats}
